% Procedencia íntegra: HIBRIDACION_ACCION_MASA_PROPAGACION.md, §§2--7.
% La gestión, el historial y los localizadores se conservan en gestion/.
\section{Action, masse et propagation : compatibilité structurelle et transport}
\label{md:hib:seccion}

\subsection{Section d’action partagée et loi des masses}
\label{md:hib:accion-masas}

Soient
\begin{equation}
 s_0=10^{-34}\mathcal U_S,\qquad
 \hbar_{\mathrm{ret}}=\eta_{\mathrm{ret}}s_0,\qquad
 R_{\mathrm{act}}=\frac{H_5}{\eta_{\mathrm{ret}}}.
 \label{md:hib:accion}
\end{equation}
Les coordonnées construites en degrés sont maintenues :
\begin{equation}
 A_{\deg}=1000\alpha,\qquad
 C^*_{\deg}=2(\eta_{\mathrm{ret}}+\alpha).
 \label{md:hib:angulos}
\end{equation}
Ce n’est qu’ensuite qu’elles sont transportées en radians au moyen de
$x=\pi A_{\deg}/180$ et $y=\pi C^*_{\deg}/180$. Les canaux sont
$q_\pm=\exp(-x\mp y)$, avec $x>|y|$. La forme elliptique s’exprime par
\begin{equation}
 \zeta=\operatorname{artanh}(y/x).
 \label{md:hib:forma}
\end{equation}
La composition utilise d’abord le coefficient adimensionnel
$\eta_{\mathrm{ret}}$ et conserve ainsi les types d’action et d’angle.
Dans cette notation, $\eta_{\mathrm{ret}}$ est le coefficient d’action,
$\eta_h(\Gamma)$ appartient à la tour d’une route et $\zeta$ exprime la
déformation de l’ellipse ; leurs fonctions restent séparées.

On travaille avec $t_0,c_{\mathrm{int}},\eta_{\mathrm{ret}}>0$, un caractère
convergent et un facteur $\Xi_\Gamma>0$. Le lecteur massique construit détermine
\begin{equation}
 \omega_0=\frac{2\pi}{108t_0},\qquad
 E_{\mathrm{clk}}=\hbar_{\mathrm{ret}}\omega_0,\qquad
 m_{\mathrm{clk}}=\frac{E_{\mathrm{clk}}}{c_{\mathrm{int}}^2}.
 \label{md:hib:reloj}
\end{equation}
Pour une route à réalisation scalaire,
\begin{equation}
 m_\Gamma^\beta=m_{\mathrm{clk}}\Xi_\Gamma,\qquad
 \Xi_\Gamma=b_e\,
 \chi_{\mathrm{ref}}(\sigma_\Gamma-\sigma_e,\eta_\Gamma-\eta_e)
 R_{\mathrm{act}}^{\kappa_\Gamma}.
 \label{md:hib:factor-ruta}
\end{equation}
$\Xi_\Gamma$ abrège le facteur structurel de route. On conserve expressément
\begin{equation}
 b_e=\frac{\varepsilon_e^{(0)}\mathcal U_E}{E_{\mathrm{clk}}}.
 \label{md:hib:normalizacion-electron}
\end{equation}
L’état central conserve sa normalisation propre par rapport au quantum
chronologique :
\begin{equation}
 \varepsilon_e^{(0)}=\frac{\sqrt3}{4}
 \left[\exp\!\left(\frac{\varphi}{\pi^2}\right)-22\alpha^3\right]
 (1+15\Delta_4),\qquad
 \varepsilon_e^\beta=\varepsilon_e^{(0)}
 R_{\mathrm{act}}^{80-54\alpha+6\Delta_4}.
 \label{md:hib:electron}
\end{equation}
Le caractère raffiné inclut les termes
\begin{equation}
 n_Ax+\frac{s_C}{6}y+k_{\Delta_4}\Delta_4+b_\zeta\zeta_{270}
 +\sum_h N_h(q_-^h-q_+^h),\qquad
 N_h=\eta_h+b_{120}\mathbf 1_{\{h=120\}},
 \label{md:hib:caracter}
\end{equation}
avec $h$ dans le semi-groupe engendré par $90$ et $120$. Le terme de hauteur
$120$ est inclus une fois. Les coefficients primitifs de retour sont
construits au moyen de
\begin{equation}
 \begin{aligned}
 a_n(\Gamma)&=\operatorname{Tr}(R_\Gamma U_\Gamma^n),\\
 p_n(\Gamma)&=\frac1n\sum_{d\mid n}\mu_{\mathrm{Mob}}(d)
                         a_{n/d}(\Gamma),\\
 \eta_h(\Gamma)&=p_h(\Gamma)-p_h(\Gamma_e),
 \end{aligned}
 \label{md:hib:retornos-primitivos}
\end{equation}
où $\mu_{\mathrm{Mob}}$ est la fonction de Möbius arithmétique.
L’action participe à l’échelle absolue, au rapport $R_{\mathrm{act}}$
et à $y$, qui détermine les canaux. Par conséquent, la simplification d’une
échelle dimensionnelle conserve les autres dépendances structurelles.
Lorsqu’on fait varier l’horloge, $b_e=\varepsilon_e^{(0)}\mathcal U_E/E_{\mathrm{clk}}$
varie avec elle : conserver la lecture électronique exige de transporter le
caractère complet, au lieu de fixer séparément ce quotient.
Dans la réalisation longitudinale intégrée, on utilise
$t_0=\pi_{\rm HMT}L_\alpha/(54c_{\mathrm{int}})$ de
\eqref{md:rigidez:reloj}.

\subsubsection{Fréquence et longueur de chaque section massive}
\label{md:hib:frecuencia-longitud}

En définissant $E_\Gamma=m_\Gamma c_{\mathrm{int}}^2$ comme énergie de repos,
$\Omega_\Gamma=E_\Gamma/\hbar_{\mathrm{ret}}$,
$\nu_\Gamma=\Omega_\Gamma/(2\pi)$ et
$R_{108}=c_{\mathrm{int}}/\omega_0$, on obtient
\begin{equation}
 \begin{aligned}
 E_\Gamma&=\hbar_{\mathrm{ret}}\omega_0\Xi_\Gamma,&
 \Omega_\Gamma&=\omega_0\Xi_\Gamma,\\
 \nu_\Gamma&=\frac{\Xi_\Gamma}{108t_0},&
 \bar\lambda_\Gamma&=
 \frac{\hbar_{\mathrm{ret}}}{m_\Gamma c_{\mathrm{int}}}
 =\frac{R_{108}}{\Xi_\Gamma}.
 \end{aligned}
 \label{md:hib:frecuencias}
\end{equation}
\begin{proof}
On substitue $m_\Gamma=m_{\mathrm{clk}}\Xi_\Gamma$ et l’on effectue les
simplifications dans une même coordinatisation dimensionnelle. Pour des routes positives $X,Y$,
la propriété de caractère et la normalisation commune donnent
\begin{equation}
 \frac{\Omega_Y}{\Omega_X}=\frac{m_Y}{m_X}
 =\chi_{\mathrm{ref}}(\sigma_Y-\sigma_X,\eta_Y-\eta_X)
 R_{\mathrm{act}}^{\kappa_Y-\kappa_X}.
 \label{md:hib:cocientes}
\end{equation}
\end{proof}
La simplification est dimensionnelle ; les canaux et la mémoire
relative demeurent. Le secteur nul reste en dehors du domaine de l’inverse de masse.
Dans une fibre non scalaire, on conserve l’opérateur ou sa mesure spectrale,
au lieu de choisir un représentant ponctuel arbitraire.

\subsubsection{Fréquence relevée et compatibilité des normalisations}
\label{md:hib:tres-frecuencias}

La construction chronologique produit également la fréquence relevée
d’holonomie
\begin{equation}
 \omega_j=\frac{\theta_j+2\pi w_j}{108t_0},\qquad
 \lambda_j=r_j\exp(i\theta_j),
 \label{md:hib:frecuencia-levantada}
\end{equation}
où $w_j$ conserve l’enroulement de l’histoire. Sa lecture énergétique est
\begin{equation}
 E_j=\hbar_{\mathrm{ret}}\omega_j
 \chi_{\mathrm{full}}R_{\mathrm{act}}^{\kappa_j}.
 \label{md:hib:energia-levantada}
\end{equation}
Cette fréquence, $\omega_0$ et $\Omega_\Gamma$, correspondent à des lecteurs
distincts. Si les deux expressions énergétiques représentent le même mode et
le même exposant, leur compatibilité exacte exige
\begin{equation}
 \frac{\omega_j}{\omega_0}\chi_{\mathrm{full}}
 =b_e\chi_{\mathrm{ref,rel}}.
 \label{md:hib:compatibilidad-normalizaciones}
\end{equation}
Les deux normalisations sont identifiées au moyen de cette égalité et ne sont pas
multipliées entre elles. La phase réduite $\theta_j$ ne récupère pas à elle seule
$w_j$. Ainsi, un retour de phase reste distinct de l’histoire
chronologique complète. L’enroulement conserve son constructeur ; cette distinction
n’autorise pas à le choisir librement ni à produire de l’énergie sans dynamique.

\subsection{Réciprocité spatiale et composition gravitationnelle}
\label{md:hib:gravedad}

La section gravitationnelle circulaire exprime le résultat du
théorème~\ref{md:rigidez:teorema} :
$G_a=c_{\mathrm{int}}^3L_\alpha^2/\hbar_a$ et $R_{108}=L_\alpha$.
La longueur provient du retour inscrit et précède cette
identification circulaire. En maintenant
$\hbar_a=\hbar_{\mathrm{ret}}$ et la convention réduite
$r_g=Gm/c_{\mathrm{int}}^2$, la masse précédente détermine
\begin{equation}
 \boxed{r_{g,\Gamma}=R_{108}\Xi_\Gamma,\qquad
 \bar\lambda_\Gamma=\frac{R_{108}}{\Xi_\Gamma},\qquad
 r_{g,\Gamma}\bar\lambda_\Gamma=R_{108}^2.}
 \label{md:hib:reciprocidad-radial}
\end{equation}
\begin{proof}
La substitution de $G_a$ et de $m_\Gamma$ donne
\[
 \frac{G_am_\Gamma}{c_{\mathrm{int}}^2}
 =\frac{c_{\mathrm{int}}^3R_{108}^2}{\hbar_{\mathrm{ret}}}
 \frac{\hbar_{\mathrm{ret}}\omega_0\Xi_\Gamma}
 {c_{\mathrm{int}}^4}
 =R_{108}\Xi_\Gamma.
\]
La seconde égalité provient de \eqref{md:hib:frecuencias} ; leur produit
prouve la troisième.
\end{proof}
La même coordonnée de route augmente une lecture radiale et contracte l’autre.
En logarithmes,
\begin{equation}
 \log\frac{r_{g,\Gamma}}{R_{108}}=\log\Xi_\Gamma,
 \qquad
 \log\frac{\bar\lambda_\Gamma}{R_{108}}=-\log\Xi_\Gamma.
 \label{md:hib:radios-logaritmicos}
\end{equation}
Le centre géométrique des deux lectures est fixé par l’aire $R_{108}^2$.

Dans la version opératoire bornée, soit $X$ autoadjoint et borné, avec
$X\geq\delta I$ pour un certain $\delta>0$, dans une fibre massive. Les opérateurs
\begin{equation}
 M=m_{\mathrm{clk}}X,\qquad
 \Lambda=R_{108}X^{-1},\qquad R_g=R_{108}X
 \label{md:hib:operadores-radiales}
\end{equation}
satisfont $R_g\Lambda=\Lambda R_g=R_{108}^2I$ sur tout l’espace.
L’affirmation utilise le calcul fonctionnel du même opérateur. Le noyau
nul reste en dehors de l’inverse. Le présent énoncé conserve le régime
borné uniformément positif ; l’extension non bornée exige de démontrer
séparément les domaines, les coupures compatibles et le passage à la limite.
La section~\ref{md:rigidez:variaciones} démontre en outre la conservation
des deux identités sous des variations non commutatives et une réduction
compatible de mémoire.

L’interprétation proposée caractérise le facteur massique relatif comme
paramètre de dilatation réciproque de ces deux lectures. La condition circulaire
appartient au résultat. Les longueurs ainsi reliées ne sont pas identifiées
par cette seule égalité à des rayons matériels mesurés ni à un horizon
électronique. L’identité radiale isolée est utilisée pour sa fonction dans la
composition, sans lui attribuer de priorité historique indépendante.

\subsection{Réponse constitutive, orientation et grandeur}
\label{md:hib:constitutivas}

Avec $T=\exp(-xI-yR)$, $R=P_+-P_-$ et les projecteurs de feuillet conservés,
le lecteur temporel est
\begin{equation}
 \begin{aligned}
 V&=(I-T^{90})(I-T^{120})^{-1}=r_+P_++r_-P_-,\\
 r_\pm&=f(s_\pm),\qquad
 s_\pm=\exp[-30(x\pm y)],\qquad
 f(s)=\frac{1-s^3}{1-s^4}.
 \end{aligned}
 \label{md:hib:operador-constitutivo}
\end{equation}
On récupère $\widehat\varepsilon=r_+^2$, $\widehat\mu=r_-^2$,
$\widehat Z=r_-/r_+$ et $\widehat c=(r_+r_-)^{-1}$.
Le lien constitutif de vitesse établit
$c_{\mathrm{int}}=c_{\mathrm{vac}}=\widehat c\,u_C$ dans la même coordinatisation.
La collection orientée de Catalan participe à $f$ par son moment
résolvant, en conservant la structure antérieure à la valeur extrême.

La décomposition exacte
\begin{equation}
 V=\widehat c^{-1/2}
 \exp\!\left[-\frac12(\log\widehat Z)R\right]
 \label{md:hib:descomposicion-v}
\end{equation}
distingue la dilatation commune et la répartition entre feuillets. L’échange des feuillets
conserve $\widehat c$ et inverse $\widehat Z$ ; la fonctionnelle de Barbero change
de signe. La section~\ref{md:compos:difeomorfismo} démontre que
$(\widehat c,\gamma)$ récupère $r_+,r_-$ dans la collection normalisée.
Sa restriction à l’image engendrée conserve cette récupérabilité ;
l’énoncé d’inversion ne déclare pas engendrée toute paire synthétique.

\subsubsection{Convexité et évaluation des canaux séparés}
\label{md:hib:convexidad}

\begin{proposition}
\label{md:hib:proposicion-convexidad}
Pour $x>0$ fixé et $|y|<x$, la fonction
$y\mapsto\log\widehat c(x,y)$ est paire et strictement convexe.
Par conséquent,
\begin{equation}
 \widehat c(x,y)\geq\widehat c(x,0)
 =\frac{1}{f(e^{-30x})^2},
 \label{md:hib:cota-velocidad}
\end{equation}
avec égalité exactement pour $y=0$.
\end{proposition}
\begin{proof}
Soit $g(u)=\log f(e^{-u})$, pour $u>0$. On a
\begin{equation}
 g'(u)=\frac{3}{e^{3u}-1}-\frac{4}{e^{4u}-1}>0,
 \label{md:hib:derivada-g}
\end{equation}
\begin{equation}
 g''(u)=-\frac{9}{4\sinh^2(3u/2)}
 +\frac{16}{4\sinh^2(2u)}<0.
 \label{md:hib:segunda-g}
\end{equation}
La première inégalité provient de la décroissance de
$a/(\exp(au)-1)$ sur $a>0$. Pour la seconde, la fonction
$a/\sinh(au/2)$ décroît strictement : sa dérivée a le signe de
$\sinh v-v\cosh v$, négatif pour $v>0$. Comme
\[
 \log\widehat c(x,y)=-g(30(x+y))-g(30(x-y)),
\]
sa dérivée seconde par rapport à $y$ est
\begin{equation}
 -900\bigl[g''(30(x+y))+g''(30(x-y))\bigr]>0.
 \label{md:hib:segunda-velocidad}
\end{equation}
La parité fixe l’unique minimum en zéro.
\end{proof}
On obtient également
\begin{equation}
 \begin{aligned}
 \log\widehat c(x,y)
 &=-2g(30x)-900g''(30x)y^2+O(y^4),\\
 \log\widehat Z(x,y)
 &=-60g'(30x)y+O(y^3).
 \end{aligned}
 \label{md:hib:desarrollos-canales}
\end{equation}
La vitesse perd le signe d’orientation, mais conserve une réponse
quadratique à la séparation des canaux. L’impédance conserve une information
orientée dès le premier ordre. Faire la moyenne de $x+y$ et $x-y$ avant d’appliquer le
lecteur élimine cette contribution. L’affirmation est une conséquence
démontrée de la réponse constitutive, distincte d’une nouvelle mesure.

Sur la branche positive d’action,
$\eta_{\mathrm{ret}}=(90/\pi)(y-x/500)$ ; par conséquent,
\begin{equation}
 \hbar_{\mathrm{ret}}=s_0\frac{90}{\pi}(y-x/500),\qquad
 E_k=s_0\frac{90u_C|k|}{\pi}(y-x/500)\widehat c(x,y).
 \label{md:hib:energia-angular}
\end{equation}
Pour $x$ fixé, $|k|>0$ et $x/500<y<x$, cette dernière fonction croît
strictement en $y$. C’est une propriété de la collection de réalisation
radiative homogène. Son application à une trajectoire engendrée conserve le
constructeur de $y$ ; la variation d’un paramètre de cette collection n’équivaut pas
à elle seule à une variation expérimentale d’une constante universelle.
Sous inversion d’orientation, on transporte également l’étiquette d’action,
en maintenant la branche positive correspondante.

\subsection{Action, horloge et information thermique}
\label{md:hib:termica}

La section thermique conserve
$k_B\Theta_{108}\log3=E_{\mathrm{clk}}$, avec
$\Theta_{108}=\Theta_{\mathrm{clk,ret}}$ dans la section de retour.
Nous définissons
$\beta_{108}=(k_B\Theta_{108})^{-1}$. Pour
$E_\Gamma=E_{\mathrm{clk}}\Xi_\Gamma$, on déduit exactement
\begin{equation}
 \beta_{108}E_\Gamma=(\log3)\Xi_\Gamma,
 \qquad
 \boxed{\exp(-\beta_{108}E_\Gamma)=3^{-\Xi_\Gamma}.}
 \label{md:hib:peso-termico}
\end{equation}
L’expression est un poids de Gibbs non normalisé. Pour une collection finie
d’états, ou une collection ayant une somme de partition convergente, et des
multiplicités $g_\Gamma$, on obtient
\begin{equation}
 Z_{\mathrm{part}}=\sum_\Gamma g_\Gamma3^{-\Xi_\Gamma},\qquad
 p_\Gamma=\frac{g_\Gamma3^{-\Xi_\Gamma}}{Z_{\mathrm{part}}},\qquad
 \frac{p_Y/g_Y}{p_X/g_X}=3^{-(\Xi_Y-\Xi_X)}.
 \label{md:hib:particion}
\end{equation}
On suppose une réalisation canonique à $\Theta_{108}$ et une origine énergétique
commune. Lorsqu’un potentiel chimique intervient, il est incorporé à l’énergie
correspondante. L’indice $\Gamma$ énumère des états ou des groupes physiques
distincts ; plusieurs histoires équivalentes ne produisent pas automatiquement une
multiplicité. Les entropies adimensionnelles sont
\begin{equation}
 \begin{aligned}
 H_{\mathrm{micro}}
 &=\log Z_{\mathrm{part}}+(\log3)\sum_\Gamma p_\Gamma\Xi_\Gamma,\\
 H_{\mathrm{grupos}}
 &=H_{\mathrm{micro}}-\sum_\Gamma p_\Gamma\log g_\Gamma.
 \end{aligned}
 \label{md:hib:entropias}
\end{equation}
Pour affirmer la finitude de l’entropie, il faut en outre
\begin{equation}
 \sum_\Gamma g_\Gamma\Xi_\Gamma3^{-\Xi_\Gamma}<\infty.
 \label{md:hib:condicion-entropia}
\end{equation}
La finitude de $Z_{\mathrm{part}}$ garantit à elle seule la normalisation.
Dans les collections finies, les deux conditions sont vérifiées.

La formule relie la structure massique à la fréquence propre et à la pondération
thermique dans le même cadre. Elle conserve la distinction entre énergie et chaleur,
entre coordonnée réelle $\Xi_\Gamma$ et dénombrement entier de trits, et entre une
réalisation à $\Theta_{108}$ et une affirmation sur la température de
l’univers. Pour des occupations bosoniques ou fermioniques libres par mode, avec
un potentiel chimique nul, la même substitution produit
\begin{equation}
 \overline n_\Gamma=\frac{1}{3^{\Xi_\Gamma}\mp1}.
 \label{md:hib:ocupaciones}
\end{equation}
Le signe moins correspond à Bose et le signe plus à Fermi. Le total du
groupe dégénéré est $g_\Gamma\overline n_\Gamma$. Dans Bose, on exige
$\Xi_\Gamma>0$ et le mode nul est traité séparément. Ce sont des réalisations
statistiques ultérieures appliquées à un spectre déjà construit.

\subsubsection{Compatibilité conjointe de quatre lectures}
\label{md:hib:cuatro-lecturas}

En composant les sections précédentes, sous les mêmes hypothèses,
\begin{equation}
 \boxed{\Xi_\Gamma=\frac{\Omega_\Gamma}{\omega_0}
 =\frac{r_{g,\Gamma}}{R_{108}}
 =\frac{R_{108}}{\bar\lambda_\Gamma}
 =\frac{E_\Gamma}{k_B\Theta_{108}\log3}.}
 \label{md:hib:compatibilidad-cuatro}
\end{equation}
Dans la collection circulaire, $\Xi=1$ identifie la section de fréquence
$\omega_0$, les rayons coïncidents $R_{108}$ et l’énergie $E_{\mathrm{clk}}$.
Son poids thermique non normalisé est $1/3$. C’est un point de compatibilité des
lectures. L’existence d’une espèce de l’atlas avec $\Xi=1$, ou d’une
seconde espèce engendrée pour chaque section réciproque $1/\Xi$, nécessite
son appartenance produite à l’atlas ; elle ne se déduit pas de cette égalité.

\subsection{Identité énergétique de l’ellipse et transport exact}
\label{md:hib:elipse}

Dans l’ellipse d’action,
\begin{equation}
 a_S=\sqrt{2\hbar}\,e^{\zeta/2},\qquad
 b_S=\sqrt{2\hbar}\,e^{-\zeta/2},\qquad
 Q=a_S\cos\theta,\qquad P=b_S\sin\theta,
 \label{md:hib:elipse-coordenadas}
\end{equation}
avec $\zeta=\operatorname{artanh}(C^*_{\deg}/A_{\deg})$, on a
$a_Sb_S=2\hbar$, une aire orientée
$\mathcal A(\theta)=\hbar\theta$ et $\mathcal A(2\pi)=h$.
La phase paramétrique $\theta$ et l’angle polaire d’un point de l’ellipse
conservent leurs significations géométriques respectives.

La réalisation LC déclarée utilise
\begin{equation}
 L_\zeta=\frac{Z_*}{\omega}e^{-\zeta},\qquad
 C_\zeta=\frac{1}{Z_*\omega}e^\zeta.
 \label{md:hib:lc}
\end{equation}
En coordonnées d’action $Q=\sqrt{Z_*}\,q$ et
$P=\Phi/\sqrt{Z_*}$,
\begin{equation}
 H_0=\frac{\omega}{2}
 \left(e^{-\zeta}Q^2+e^\zeta P^2\right),\qquad
 U_E=\hbar\omega\cos^2\theta,\qquad
 U_B=\hbar\omega\sin^2\theta.
 \label{md:hib:energia-elipse}
\end{equation}
On obtient $H_0=\hbar\omega$ par composition des deux contributions.
L’horloge et le transducteur appartiennent à la réalisation déclarée ; la section
$\hbar$ conserve sa construction antérieure au circuit. L’impédance
$Z_{\mathrm{LC}}/Z_*=e^{-\zeta}$ et l’impédance constitutive du vide
$f(s_-)/f(s_+)$ sont des lecteurs distincts. Leur antécédent angulaire commun
n’établit pas d’égalité entre les deux.

\subsubsection{Transport discret de la forme quadratique}
\label{md:hib:transporte-discreto}

\begin{proposition}
\label{md:hib:proposicion-discreta}
Soient $S_\zeta=\operatorname{diag}(e^{\zeta/2},e^{-\zeta/2})$,
$g_\zeta=S_\zeta^{-T}S_\zeta^{-1}$ et
$\mathsf R(-\theta)$ une rotation. Pour des valeurs de forme $\zeta_n$ et des
phases $\theta_n$, le transport
\begin{equation}
 X_{n+1}=S_{\zeta_{n+1}}\mathsf R(-\theta_n)S_{\zeta_n}^{-1}X_n
 =T_nX_n
 \label{md:hib:transporte}
\end{equation}
conserve exactement
$I_n=X_n^Tg_{\zeta_n}X_n/2$ et a pour déterminant un.
\end{proposition}
\begin{proof}
La substitution et l’identité $\mathsf R^T\mathsf R=I$ donnent
\begin{equation}
 T_n^Tg_{\zeta_{n+1}}T_n=g_{\zeta_n}.
 \label{md:hib:isometria-transportada}
\end{equation}
La forme quadratique est conservée lorsqu’on évalue les deux membres sur $X_n$.
Les facteurs de $T_n$ ont pour déterminant un, d’où
$\det T_n=1$.
\end{proof}
C’est une composition du transport radial préalablement construit. Si
$\zeta_n,\theta_n$ sont extraits d’une route TPK, l’identité s’applique à
cette séquence ; la proposition algébrique maintient séparée la sélection
de la route.

\subsubsection{Transport différentiable et terme de connexion}
\label{md:hib:conexion}

Si $\zeta(t)$ est de classe $C^1$ et $\omega(t)>0$ est continue, la relation
$X=S_\zeta Y$ et l’évolution $\dot Y=-\omega J_0Y$, avec
\begin{equation}
 J_0=\begin{pmatrix}0&-1\\1&0\end{pmatrix},
 \label{md:hib:j-canonico}
\end{equation}
donnent
\begin{equation}
 \dot Q=\frac{\dot\zeta}{2}Q+\omega e^\zeta P,\qquad
 \dot P=-\omega e^{-\zeta}Q-\frac{\dot\zeta}{2}P.
 \label{md:hib:ecuaciones-transporte}
\end{equation}
Avec les conventions $\dot Q=\partial_PH$ et
$\dot P=-\partial_QH$, l’hamiltonien est, à un terme dépendant
seulement du temps près,
\begin{equation}
 \boxed{H_{\mathrm{tot}}=
 \frac{\omega}{2}\left(e^{-\zeta}Q^2+e^\zeta P^2\right)
 +\frac{\dot\zeta}{2}QP.}
 \label{md:hib:hamiltoniano-conexion}
\end{equation}
L’invariant
\begin{equation}
 I=\frac12\left(e^{-\zeta}Q^2+e^\zeta P^2\right)
 \label{md:hib:invariante}
\end{equation}
est conservé exactement, sans nécessiter d’approximation adiabatique.
\begin{proof}
En dérivant, on obtient
\[
 \dot I=\frac{\dot\zeta}{2}
 \left(e^\zeta P^2-e^{-\zeta}Q^2\right)
 +e^{-\zeta}Q\dot Q+e^\zeta P\dot P.
\]
La substitution de \eqref{md:hib:ecuaciones-transporte} annule les termes
proportionnels à $\dot\zeta$ et les deux termes $\omega QP$.
Ainsi, $\dot I=0$. Si l’on omet le terme mixte de
\eqref{md:hib:hamiltoniano-conexion}, il reste
\begin{equation}
 \dot I=\frac{\dot\zeta}{2}
 \left(e^\zeta P^2-e^{-\zeta}Q^2\right).
 \label{md:hib:defecto-sin-conexion}
\end{equation}
Le signe mixte contraire double ce défaut.
\end{proof}
Sur l’ellipse, $QP=I\sin(2\theta)$ et $\dot\theta=-\omega$, de sorte que
\begin{equation}
 H_{\mathrm{tot}}=I\omega+\frac{I\dot\zeta}{2}\sin(2\theta).
 \label{md:hib:hamiltoniano-elipse}
\end{equation}
La matrice de la forme quadratique a pour déterminant
$\omega^2-\dot\zeta^2/4$. Pour $\omega>0$, elle est définie positive
exactement lorsque $|\dot\zeta|<2\omega$, et semi-définie positive dans le cas
d’égalité. Ce seuil concerne la positivité : l’existence du
transport et la stabilité physique conservent leurs conditions respectives.
La conservation de $I$ n’implique pas la conservation de $H_{\mathrm{tot}}$,
qui dépend explicitement du temps.

La version quantique exige la symétrisation du terme mixte,
$\dot\zeta(QP+PQ)/4$. Cette expression indique l’opération requise sans
présupposer de représentation quantique supplémentaire.

\subsubsection{Compatibilité avec la section d’action génératrice du contre-angle}
\label{md:hib:compatibilidad-accion}

Dans la construction directrice,
\begin{equation}
 \zeta=\operatorname{artanh}
 \left[\frac{2(\eta_{\mathrm{ret}}+\alpha)}{1000\alpha}\right].
 \label{md:hib:forma-genealogica}
\end{equation}
Fixer $\alpha$ et $\eta_{\mathrm{ret}}$ fixe également $\zeta$.
Le résultat avec $\zeta(t)$ variable est donc une extension de
transport entre formes. Pour représenter une trajectoire physique de la paire
directrice, la variation doit provenir de la généalogie correspondante, et non
d’une déformation choisie de manière indépendante.

Le transport peut également être composé exactement lorsque la
section positive d’action change entre deux états, en conservant la même coordinatisation
dimensionnelle. Soient
\begin{equation}
 \hbar_n>0,\qquad
 \rho_{\mathrm{act},n}=\frac{\hbar_{n+1}}{\hbar_n},\qquad
 \widetilde T_n=\sqrt{\rho_{\mathrm{act},n}}\,
 S_{\zeta_{n+1}}\mathsf R(-\theta_n)S_{\zeta_n}^{-1}.
 \label{md:hib:transporte-conforme}
\end{equation}
Alors
\begin{equation}
 \det\widetilde T_n=\rho_{\mathrm{act},n},\qquad
 \widetilde T_n^Tg_{\zeta_{n+1}}\widetilde T_n
 =\rho_{\mathrm{act},n}g_{\zeta_n},\qquad
 \frac{I_{n+1}}{\hbar_{n+1}}=\frac{I_n}{\hbar_n}.
 \label{md:hib:accion-normalizada}
\end{equation}
Les égalités découlent de \eqref{md:hib:isometria-transportada} en multipliant
par le facteur scalaire. C’est un transport conformément symplectique : il conserve
l’action normalisée et conserve l’aire d’action non normalisée
exactement lorsque $\rho_{\mathrm{act},n}=1$.
Sa limite différentiable a pour trace $\dot\hbar/\hbar$. Un générateur
hamiltonien linéaire sur un plan à forme symplectique fixe a une trace nulle ;
ainsi, une variation matérielle de $\hbar$ exige de spécifier la forme
transportée ou les degrés de liberté avec lesquels l’action est échangée.
Le terme mixte de déformation conserve à lui seul une trace nulle.

Le résultat est un critère de compatibilité des réalisations. Pour
interpréter des variations numériques dues à un changement d’unités, on
transporte d’abord les bases : le changement de coordinatisation maintient la section
physique. La composition n’affirme pas une variabilité expérimentale de la
constante de Planck.

\subsubsection{Transformation de représentation et déformation matérielle}
\label{md:hib:interpretacion-transporte}

Si $S_\zeta$ est un changement passif de coordonnées, le terme mixte est une
connexion de représentation. Les observables sont également transportées,
\begin{equation}
 L_X(S_\zeta Y)=L_Y(Y).
 \label{md:hib:observables-transportados}
\end{equation}
Cette égalité conserve la lecture physique lorsqu’on change l’expression de
l’hamiltonien, sans introduire d’interaction supplémentaire ni de travail physique.

Une déformation matérielle exige d’obtenir $\zeta$ à partir de la mise à jour TPK et
de montrer un changement relatif entre système et référence, tous deux transportés
de manière cohérente. On peut confronter une phase relative, une réponse
$Z_{\mathrm{sistema}}/Z_{\mathrm{referencia}}$ ou un transfert entre
modes définis par le détecteur. Il faut distinguer deux histoires TPK non
équivalentes par changement de représentation et conserver leurs conditions de
lecture. Lorsqu’un échange énergétique est attribué, le bilan comprend
champ, mémoire et dispositif. Le théorème de transport délimite cette
distinction et maintient séparée la loi dynamique qui la rend physique.

\subsection{Quantité de mouvement, confrontations et portée de la composition}
\label{md:hib:momento}

La réalisation de Clifford--Dirac reçoit la masse de route. Dans une même
section d’action, son coefficient s’écrit
\begin{equation}
 \frac{m_\Gamma c}{\hbar}
 =\frac{\omega_0\Xi_\Gamma}{c}.
 \label{md:hib:coeficiente-dirac}
\end{equation}
Dans une réalisation libre et homogène, avec la relation relativiste de
dispersion déclarée, celle-ci se réexprime comme
\begin{equation}
 \Omega^2=c^2|k|^2+\omega_0^2\Xi_\Gamma^2.
 \label{md:hib:dispersion}
\end{equation}
La réécriture conserve les hypothèses de cette réalisation et ses antécédents
de connexion spatio-temporelle. Les moments de retour $q_-^h-q_+^h$
sont des évaluations des tours et se distinguent des moments mécaniques.

Le facteur $\Xi_\Gamma$ réapparaît comme poids massique, fréquence relative,
dilatation radiale et exposant thermique. La réponse constitutive détermine
la conversion spatiale commune. Barbero conserve une information orientée
complémentaire à la vitesse. L’action fixe la mesure aréale ; l’horloge
introduit le taux temporel ; la forme et son transport spécifient la répartition
et son évolution.

Les confrontations correspondantes comprennent les quotients de masses et de fréquences ;
la compatibilité de la vitesse, de l’impédance et de Barbero ; la conservation de la forme
quadratique sous transport ; et les poids thermiques une fois la température et les
multiplicités fixées. Ces relations préservent plus de données qu’une coïncidence
décimale isolée. Leurs preuves maintiennent le bilan énergétique et les
conditions des entropies considérées ; elles n’impliquent pas une entropie
thermodynamique nulle.
